\section{From $Y^+(\widehat{\mathfrak{gl}}_1)$ to a chiral Yangian}
\label{wn:sec:yplus-to-chiral-yangian-treatise}

\providecommand{\kBKM}{\kappa_{\mathrm{BKM}}}
\providecommand{\kch}{\kappa_{\mathrm{ch}}}
\providecommand{\kcat}{\kappa_{\mathrm{cat}}}
\providecommand{\kfib}{\kappa_{\mathrm{fiber}}}
\providecommand{\gghone}{\widehat{\mathfrak{gl}}_1}
\providecommand{\Yplus}{Y^+_{\epsilon_1, \epsilon_2}}
\providecommand{\Yfull}{Y_{\epsilon_1, \epsilon_2}}
\providecommand{\Sh}{\mathrm{Sh}}
\providecommand{\Hilb}{\mathrm{Hilb}}
\providecommand{\CoHA}{\mathrm{CoHA}}
\providecommand{\Fact}{\mathrm{Fact}}
\providecommand{\ChirY}{\mathcal{Y}^{\mathrm{ch}}}

The passage from the positive half of the affine Yangian to a chiral
Yangian has three separate operations: Drinfeld doubling, choice of
completion and Cartan data, and a vertex or defect realisation. The
positive half alone is not a vertex algebra and is not
$\mathcal W_{1+\infty}$.

\subsection{Equivariant parameters and the CY$_3$ constraint}
\label{wn:subsec:yplus-setup}

Let $\mathbb{F}:=\mathbb{C}(\epsilon_1,\epsilon_2)$ and
$\epsilon_3:=-\epsilon_1-\epsilon_2$. The relation
$\epsilon_1+\epsilon_2+\epsilon_3=0$ is the Calabi--Yau threefold
condition: it is the weight-space relation for the anti-canonical
subtorus
\[
T^2=\{t\in T^3:t_1t_2t_3=1\}
\]
acting on $\mathbb C^3$.

\emph{Shuffle kernel.}
\[
\omega(z,w):=
\frac{(z-w-\epsilon_1)(z-w-\epsilon_2)(z-w-\epsilon_3)}{(z-w)^3}.
\]

\emph{Shuffle algebra.}
Let $\Lambda_n:=\mathbb F[z_1,\ldots,z_n]^{S_n}$ and
\[
\Sh:=\bigoplus_{n\ge0}\Lambda_n.
\]
For $F\in\Lambda_m$ and $G\in\Lambda_n$, define
\[
(F\star G)(z_1,\ldots,z_{m+n})
:=
\frac{1}{m!n!}\sum_{\sigma\in S_{m+n}}
\sigma\!\left(
F(z_1,\ldots,z_m)G(z_{m+1},\ldots,z_{m+n})
\prod_{1\le i\le m<j\le m+n}\omega(z_j,z_i)
\right).
\]
Then $(\Sh,\star)$ is an associative unital $\mathbb F$-algebra.
Associativity is the direct shuffle calculation with the kernel
$\omega$.

\emph{Positive half.}
\[
\Yplus:=\Sh.
\]
The name ``positive half of $Y(\gghone)$'' records the later theorem
identifying $\Sh$ with the upper triangular part of the Yangian in its
Drinfeld-current presentation. It is not the definition.

Different shuffle kernels in the literature differ by signs, rescaling
of the $\epsilon_i$, and the point at which the CY$_3$ relation is
imposed. The form above matches Negu\c{t} 2014
\texttt{arXiv:1405.1989}; Feigin--Hashizume--Hoshino--Shiraishi--
Yanagida 2009 uses an equivalent rescaling.

\subsection{Drinfeld double}
\label{wn:subsec:yplus-drinfeld-double}

The Drinfeld double construction (Drinfeld 1987 ICM; Kassel 1995,
\emph{Quantum Groups}, §IX.4) takes bialgebras $A_+$, $A_-$ and a
non-degenerate Hopf pairing
\[
\langle\cdot,\cdot\rangle:A_+\otimes A_-\to k
\]
and produces a Hopf algebra $D(A_+,A_-)$ with underlying vector space
$A_+\otimes A_-$. The multiplication is
\[
(a_+\otimes a_-)(b_+\otimes b_-)
=\sum
\langle a_-^{(1)},b_+^{(1)}\rangle
\langle S(a_-^{(3)}),b_+^{(3)}\rangle
a_+b_+^{(2)}\otimes a_-^{(2)}b_-.
\]

For $\Yplus$, the required data are:
\begin{enumerate}
\item a coproduct $\Delta_+$ on $\Sh$;
\item a negative bialgebra $Y^-\cong\Sh^{\mathrm{op}}$;
\item a non-degenerate Hopf pairing
    $\langle,\rangle:\Yplus\otimes Y^-\to\mathbb F$;
\item a completion on which the pairing is non-degenerate.
\end{enumerate}

Negu\c{t} 2014 Thm.\ 3.4 and Schiffmann--Vasserot 2013
\texttt{arXiv:1202.2756} give the topological coproduct, the Hopf
pairing with the opposite half, and the Cartan extension
\[
Y^0\cong \mathbb F[\![\psi_0,\psi_1,\ldots]\!].
\]
The completed double has triangular decomposition
\[
Y(\gghone)=\Yplus\otimes Y^0\otimes Y^-.
\]
The pairing is not defined on the rational shuffle algebra without
completion. A specified $\widehat{\Sh}$, analogous to a formal Laurent
completion, is part of the object.

\subsection{Drinfeld currents}
\label{wn:subsec:yplus-loop-presentation}

For the current presentation of $Y(\gghone)$, introduce
\[
e(z)=\sum_{k\ge0}e_kz^{-k-1},\qquad
f(z)=\sum_{k\ge0}f_kz^{-k-1},\qquad
\psi(z)=\exp\!\left(\sum_{k\ge0}h_kz^{-k-1}\right).
\]
Set
\[
\varphi(u):=\frac{(u+\epsilon_1)(u+\epsilon_2)(u+\epsilon_3)}{u^3}.
\]
The relations are:
\begin{enumerate}[label=(R\arabic*)]
\item $[\psi(z),\psi(w)]=0$;
\item $\psi(z)e(w)\psi(z)^{-1}=\varphi(z-w)e(w)$ and
      $\psi(z)f(w)\psi(z)^{-1}=\varphi(z-w)^{-1}f(w)$;
\item $(z-w-\epsilon_1)(z-w-\epsilon_2)(z-w-\epsilon_3)e(z)e(w)$
      \[
      =-(w-z-\epsilon_1)(w-z-\epsilon_2)(w-z-\epsilon_3)e(w)e(z),
      \]
      and similarly for $f$;
\item $[e(z),f(w)]=\dfrac{\psi(z)-\psi(w)}{z-w}$;
\item $\mathrm{Sym}_{z_1,z_2,z_3}
      (z_1-2z_2+z_3)[e(z_1),[e(z_2),e(z_3)]]=0$,
      and similarly for $f$.
\end{enumerate}
Tsymbaliuk 2017 \texttt{arXiv:1703.04551} Thm.\ 1.1 identifies this
topological Hopf algebra with the completed Drinfeld double
$\Yplus\bowtie Y^0\bowtie Y^-$.

\subsection{Vacuum module and vertex algebra}
\label{wn:subsec:yplus-vacuum-vs-vertex}

A module of $Y(\gghone)$ is not a vertex algebra. Fix the trivial
character $\chi_0:Y^0\to\mathbb C$, $\psi_k\mapsto0$. The Verma module
\[
M(\chi_0):=Y(\gghone)\otimes_{Y^0\otimes Y^-}\mathbb C_{\chi_0}
\]
is the vacuum Verma; its irreducible quotient is $L(\chi_0)$.

A vertex algebra is a vector space $V$ with a vacuum vector
$|0\rangle$, a state-to-field map
$Y(\cdot,z):V\to\mathrm{End}(V)[[z,z^{-1}]]$, a translation operator,
and the vacuum, translation, and locality axioms. To turn a Yangian
module into a vertex algebra one must define the fields and prove
locality. This is extra structure.

\subsection{Hilbert-scheme action}
\label{wn:subsec:yplus-MO-action}

Let
\[
V:=\bigoplus_{n\ge0}H^*_T(\Hilb^n(\mathbb C^2),\mathbb F),
\]
with $T$ acting on $\mathbb C^2$ with weights
$\epsilon_1,\epsilon_2$. Schiffmann--Vasserot 2012, Maulik--Okounkov
2012 §14 give an action
\[
\rho:Y(\gghone)\to\mathrm{End}_{\mathbb F}(V)
\]
by explicit correspondences on
$\Hilb^n\times\Hilb^{n+1}$.

Independently, Nakajima 1997 §9 and Grojnowski 1996 give a Heisenberg
vertex algebra structure on $V$:
\[
[\alpha_m,\alpha_n]=m\delta_{m,-n}.
\]
Thus $V$ is both a $Y(\gghone)$-module and the free-boson Fock VOA.
The Yangian acts by operators compatible with this VOA structure; the
Yangian itself is not thereby a vertex algebra.

\subsection{Gauge-theoretic chiral Yangian}
\label{wn:subsec:yplus-costello}

Costello 2013 \texttt{arXiv:1303.2632}, Costello--Gwilliam 2017, and
the $\Omega$-deformed upgrades of Costello--Gaiotto 2018 and
Costello--Paquette 2020 construct a factorisation algebra
\[
\mathcal F^{U(1),\Omega}_{4D}
\]
on the spectral line $\mathbb C_z$ from the holomorphic-topological
twist of $4D$ $\mathcal N=2$ $U(1)$ gauge theory on
$\mathbb C^2_{\epsilon_1,\epsilon_2}$. Costello 2013 §11 gives an
action of $Y(\gghone)$ on this factorisation algebra, with
$\epsilon_3=-\epsilon_1-\epsilon_2+\hbar$ in the quantum deformation.

A holomorphic factorisation algebra on $\mathbb C$ satisfying locality
is equivalent to a vertex algebra (Costello--Gwilliam 2017 Vol.\ I
Thm.\ 5.3.3). Define
\[
\ChirY^{\mathrm{Cost}}
\]
as the resulting vertex algebra. Its identification with
$\mathcal W_{1+\infty}$ requires the defect construction and the
state-field realisation; it is not a consequence of
$\CoHA(\mathbb C^3)\simeq Y^+$ alone.

\subsection{Kapranov--Vasserot route}
\label{wn:subsec:yplus-KV}

Kapranov--Vasserot 2019 \texttt{arXiv:1901.07641} constructs the CoHA
of $\mathbb C^3$ by factorisation methods. This supplies additional
factorisation structure on the CoHA package; it does not change the
Vol III specialised output rule that a CY$_3$ chiral image is $\Eone$.
The $\Etwo$ structure relevant to the full chiral Yangian is recovered
on the centre, double, or defect category under the stated hypotheses.

The comparison
\[
\int_{\mathbb C}\CoHA(\mathbb C^3)\cong\ChirY^{\mathrm{KV}}
\]
is a proof obligation at the level of factorisation algebras. The
comparison
\[
\ChirY^{\mathrm{Cost}}\cong\ChirY^{\mathrm{KV}}
\]
is another proof obligation in the Costello--Paquette framework.

\subsection{Gaiotto--Rap\v{c}\'ak $Y$-algebras}
\label{wn:subsec:yplus-GR}

Gaiotto--Rap\v{c}\'ak 2017 \texttt{arXiv:1703.00982} defines
vertex algebras
\[
Y_{L,M,N}^{\epsilon_1,\epsilon_2,\epsilon_3}
\]
as local observables of $4D$ $\mathcal N=2$ gauge theory along a
Y-shaped divisor in $\mathbb C^3$ with ramification data $(L,M,N)$.
The identification
\[
Y_{0,0,0}^{\epsilon_1,\epsilon_2,\epsilon_3}
\cong
\ChirY^{\mathrm{Cost}}
\]
is supported by low-order OPE computations. A complete theorem requires
the full matching of generators, OPEs, central charge normalisation, and
factorisation products.

\subsection{Identification sources and proof obligations}

\begin{center}
\begin{tabular}{l|l|l}
\textbf{Identification} & \textbf{Mathematical scope} & \textbf{Source} \\\hline
$\Yplus\cong\Sh$ & definition plus shuffle calculation & Negu\c{t} 2014 \\
$Y(\gghone)\cong D(\Yplus)$ & theorem after completion & SV 2013 \\
$Y(\gghone)$ current presentation & theorem & Tsymbaliuk 2017 \\
$Y(\gghone)\curvearrowright\bigoplus H^*_T(\Hilb^n)$ & theorem & SV 2012, MO 2012 \\
$V$ is the Heisenberg VOA & theorem & Nakajima 1997, Grojnowski 1996 \\
$\mathcal F^{U(1),\Omega}_{4D}$ is a factorisation algebra & theorem & Costello 2013, CG 2017 \\
$\mathcal F^{U(1),\Omega}_{4D}$ is a vertex algebra & theorem via holomorphic locality & CG 2017 Vol.\ I \\
$Y(\gghone)$ deforms $\mathcal F^{U(1),\Omega}_{4D}$ & theorem & Costello 2013 §11 \\
$\ChirY^{\mathrm{Cost}}\cong\ChirY^{\mathrm{KV}}$ & proof obligation & Costello--Paquette 2020 \\
$Y_{0,0,0}^{\mathrm{GR}}\cong\ChirY^{\mathrm{Cost}}$ & proof obligation & GR 2017 framework \\
$V\cong\ChirY^{\mathrm{Cost}}$ directly & proof obligation & low-order checks only \\
\end{tabular}
\end{center}

\subsection{Compact $K3\times E$ comparison}

For the compact Calabi--Yau threefold $K3\times E$, the analogy with
$\mathbb C^3$ has the following form.
\begin{itemize}
\item $\CoHA(\mathbb C^3)$ has a compact analogue
    $\CoHA(K3\times E)$ only after the proper critical-cohomology
    hypotheses of Davison's framework are supplied.
\item $Y^+(\gghone)$ suggests a positive-half BPS algebra for
    $K3\times E$; this remains a construction problem.
\item A Drinfeld double for the compact case requires a Hopf pairing
    from Serre duality on the moduli stacks.
\item The chiral or vertex algebra at the $\Delta_5$ locus requires a
    defect/factorisation construction. The BKM Lie algebra
    $\mathfrak g_{\Delta_5}$ is not itself the chiral algebra.
\end{itemize}

The proved compact statement is numerical: Oberdieck--Pixton 2017 gives
the reduced DT series, and the Borcherds product for
$\Phi_{10}=\Delta_5^2$ gives root multiplicities. This identifies
graded dimensions. A Lie-bracket isomorphism
\[
\mathfrak g_{\mathrm{BPS}}(K3\times E)\cong\mathfrak g_{\Delta_5}
\]
requires a bracket comparison for
\[
[\mathfrak g_{\mathrm{BPS},\gamma_1},
 \mathfrak g_{\mathrm{BPS},\gamma_2}]
\to
\mathfrak g_{\mathrm{BPS},\gamma_1+\gamma_2}
\]
against the BKM bracket. A chiral-BKM statement further requires
fields, OPEs, locality, and the factorisation algebra on a chosen
spectral curve.
